% Source: https://artofproblemsolving.com/community/c3h1924330_simple_angle_chasing_problem
% Stable identity: aops-topic-1924330-q01
Since $\angle B$ is acute, the inscribed angle $\angle ABC$ and the central angle $\angle AOC$ stand on the same arc $AC$ (the one not containing $B$), so $\angle AOC=2\angle ABC<180^\circ$. Triangle $AOC$ is isosceles with $OA=OC$, so
\[
\angle CAO=\dfrac{180^\circ-2\angle B}{2}=90^\circ-\angle B .
\]

\begin{asy}
usepackage("amsmath");
size(6cm);
defaultpen(fontsize(10pt));
real R=3;pair O=(0,0);
pair A=R*dir(105),B=R*dir(215),C=R*dir(-35);
pair D=B+dot(A-B,C-B)/dot(C-B,C-B)*(C-B);
draw(circle(O,R),gray(0.5));
draw(A--B--C--cycle,linewidth(1));
draw(A--D,linewidth(0.8));draw(A--O,linewidth(0.8));draw(O--C,dashed);
pair u=unit(C-B),v=unit(A-D);
draw(D+0.3*u--D+0.3*u+0.3*v--D+0.3*v);
dot(O);dot(D);
label("$A$",A,N);label("$B$",B,SW);label("$C$",C,SE);label("$O$",O,E);label("$D$",D,S);
label("\mbox{}",(min(currentpicture,user=true)+max(currentpicture,user=true))/2,black); // keeps the text after the figure black when LaTeX typesets the labels inline
\end{asy}

Hence $\angle BAD=90^\circ-\angle B$. In triangle $ABD$ the angle at $B$ is $\angle ABC$ (because $D$ lies on segment $BC$), so
\[
\angle ADB=180^\circ-\angle B-(90^\circ-\angle B)=90^\circ ,
\]
that is, $AD\perp BC$. (The hypothesis on $\angle C$ guarantees that the foot of the altitude from $A$, the point $D$, really lies on segment $BC$.)
